% --- UNIVERSAL PREAMBLE BLOCK ---
\documentclass[11pt, a4paper]{article}
\usepackage[a4paper, top=2.5cm, bottom=2.5cm, left=2cm, right=2cm]{geometry}
\usepackage{fontspec}

\usepackage[bidi=bidi, provide=*]{babel}
\babelprovide[import, main]{english}

% Set default/Latin font to Sans Serif in the main (rm) slot
\babelfont{rm}{Noto Sans}

% Required packages for tables, math, and layout
\usepackage{amsmath}
\usepackage{amssymb}
\usepackage{booktabs}
\usepackage{tabularx}
\usepackage{hyperref}
\usepackage{enumitem}
\usepackage{parskip}

% Setup hyperref
\hypersetup{
    colorlinks=true,
    linkcolor=blue,
    filecolor=magenta,
    urlcolor=blue,
    pdftitle={Code-Based Material Science Project Plan},
}

\begin{document}

\title{\textbf{Computational Materials Science and Simulation Architectures for Cryogenic Quantum Metamaterials}}
\date{}
\maketitle

\noindent\textbf{Project disclaimer.} This document is research/vision material. The QASIC Engineering-as-Code project is a solo-developer, digital-twin effort with no lab, company, or partnerships; references to facilities, vendors, or companies are for research context only and do not imply affiliation or partnership.

\section{The Paradigm Shift to Code-Based Material Science in Quantum Engineering}

The realization of fault-tolerant quantum computing and global quantum communication networks is currently constrained by two distinct yet mathematically parallel physical bottlenecks: the micro-scale physical limitations of scaling solid-state quantum processors and the macro-scale environmental limitations of distributing quantum entanglement over global distances.\textsuperscript{1} In the computational domain, the transition from intermediate-scale noisy quantum systems to fully fault-tolerant architectures---which inherently require the precise control and entanglement of millions of physical qubits---is actively hindered by a severe physical wiring crisis.\textsuperscript{1} Driving modern superconducting qubits currently mandates a brute-force connectivity approach, utilizing numerous physical coaxial interconnects routed from room-temperature control electronics down to the deepest millikelvin stages of a dilution refrigerator.\textsuperscript{1}

This static wiring paradigm dissipates active power and conducts passive heat, creating an insurmountable thermodynamic burden. Even when utilizing highly specialized low-thermal-conductivity alloys, a single UT-085-SS-SS coaxial cable can introduce nearly 1 milliwatt (mW) of heat load to the 4 Kelvin (K) stage of a cryogenic system.\textsuperscript{1} When this heat load is multiplied by the thousands of interconnects required for comprehensive error correction, it rapidly overwhelms the micro-watt cooling capacities available at the lowest cryogenic stages.\textsuperscript{1} Furthermore, the physical footprint of these interconnects restricts multi-qubit gates to static, permanently etched planar topologies, severely limiting nearest-neighbor interactions and preventing dynamic, all-to-all qubit connectivity.\textsuperscript{1}

Concurrently, the distribution of continuous-variable quantum states over macroscopic distances via Satellite Communications (SATCOM) is severely hindered by atmospheric Mie scattering, an optical phenomenon that occurs because the operational wavelengths of standard optical Quantum Key Distribution (QKD) photons are roughly equivalent to the physical size of atmospheric water droplets and suspended aerosols.\textsuperscript{1} While shifting to microwave photons allows the signals to effectively diffract around these suspended particles due to the Rayleigh Scattering Approximation, their exceptionally low energy profile renders them highly susceptible to ambient thermal noise, traditionally preventing their use in over-the-air quantum communication.\textsuperscript{1}

To decouple physical scaling from these dual thermodynamic and atmospheric overloads, advanced quantum engineering is aggressively shifting toward a ``Quantum Application-Specific Integrated Circuit'' (Quantum ASIC) paradigm.\textsuperscript{1} By entirely replacing static physical waveguides with a programmable holographic metasurface functioning as a software-defined quantum bus, systems can dynamically manipulate the local density of photonic and phononic states at the sub-wavelength scale.\textsuperscript{1} However, the physical prototyping of such deep-cryogenic architectures---relying on radio-frequency Superconducting Quantum Interference Devices (rf-SQUIDs), piezoelectric lithium niobate ($\text{LiNbO}_3$) Bulk Acoustic Wave (BAW) resonators, and 28nm Fully Depleted Silicon-On-Insulator (FD-SOI) Cryo-CMOS controllers---is prohibitively slow and highly resource-intensive.\textsuperscript{1}

The extreme conditions of a 10 millikelvin (mK) operational environment cause traditional control materials, such as nematic liquid crystals or phase-change vanadium dioxide, to freeze completely or exhibit anomalous defect behaviors, rendering conventional macroscopic electromagnetics intuition highly inaccurate.\textsuperscript{1} Consequently, the advancement of the Quantum ASIC strictly requires a massive pivot toward code-based materials science. High-performance computational frameworks, open-source Python libraries, and Deep Neural Network (DNN) surrogate models are now the primary vehicles for characterizing cryogenic metamaterials, mapping algorithmic topologies, and synthesizing sub-wavelength phase profiles.\textsuperscript{1} This exhaustive research report investigates the computational physics, numerical solvers, and software engineering frameworks that drive the \textit{in-silico} realization of cryogenic holographic metasurfaces and quantum SATCOM apertures.

\section{The Quantum ASIC Paradigm: Protocol Logic and Restricted Topologies}

Before simulating the underlying quantum materials, the architecture of the Quantum ASIC must be programmatically defined at the logical level. The Quantum ASIC is explicitly designed to simplify quantum hardware to the absolute minimum required for specific cryptographic and routing protocols, abandoning the overhead of universal quantum computation in favor of highly optimized, task-specific execution.\textsuperscript{1} In code-based material science, this logical layer is defined through abstract Python modules that dictate the allowed interactions before any electromagnetic simulation occurs.

The foundational hardware topology is abstracted in simulation modules such as \texttt{asic/topology.py} and \texttt{asic/gate\_set.py}.\textsuperscript{1} These software constructs define a minimal physical resource specification: exactly three logical qubits arranged in a strict linear chain configuration represented as $0 - 1 - 2$.\textsuperscript{1} Within this restricted topology, only adjacent pairs of physical qubits are capable of native interaction.\textsuperscript{1} To enforce hardware limitations during simulation, the software intentionally excludes universal gate libraries, restricting the 1-qubit operations to Hadamard (H), Pauli-X (X), Pauli-Z (Z), and parameterized rotations around the X-axis ($R_x(\theta)$).\textsuperscript{1} Two-qubit operations are strictly limited to Controlled-NOT (CNOT) gates applied exclusively to adjacent edges, specifically the $(0,1)$ or $(1,2)$ qubit pairs.\textsuperscript{1} All-to-all connectivity, SWAP gates, T gates, and Toffoli gates are programmatically excluded, forcing the metasurface routing layer to bridge non-adjacent interactions via holographic projection rather than hardcoded physical microwave traces.\textsuperscript{1}

The utility of this abstracted code layer is its ability to validate quantum protocols before they are mapped to the physical metamaterial. Modules such as \texttt{asic/circuit.py} and \texttt{asic/executor.py} instantiate \texttt{Op} dataclasses to execute step-by-step state vector simulations.\textsuperscript{1} Through these programmatic structures, the system validates the precise sequences required for specialized secure communications.

\begin{table}[htbp]
\centering
\renewcommand{\arraystretch}{1.5}
\begin{tabularx}{\textwidth}{>{\raggedright\arraybackslash}p{3cm} >{\raggedright\arraybackslash}X >{\raggedright\arraybackslash}X >{\raggedright\arraybackslash}X}
\toprule
\textbf{Protocol Designation} & \textbf{Qubit Role Assignment} & \textbf{Sequential ASIC Gate Operations} & \textbf{Core Application} \\
\midrule
Quantum Teleportation & Qubit 0 (Message), Qubit 1 (Alice's Bell half), Qubit 2 (Bob's Bell half) & H(1), CNOT(1,2), CNOT(0,1), H(0) & Establishing continuous-variable entanglement across the routing bus.\textsuperscript{1} \\
Quantum Bit Commitment & Qubit 0 (Alice's committer), Qubit 1 (Bob's receiver) & H(0), CNOT(0,1) (Optional X(0) to commit-to-1) & Secure cryptographic handshakes utilizing the minimal topology.\textsuperscript{1} \\
Tamper-Evidence (Thief) & Qubits 0, 1, and 2 identical to Teleportation roles & Teleport ops + $R_x(\theta)$ applied to Qubit 2 & Detecting eavesdropping via fidelity drops induced by channel disturbance.\textsuperscript{1} \\
\bottomrule
\end{tabularx}
\caption{Protocol Validations within the Software-Defined Quantum ASIC Layer}
\end{table}

The logic governing these operations relies heavily on state tracking matrices. For instance, in the ``Thief'' tamper-evidence protocol script (\texttt{tamper\_evident.py}), the software simulates an interception of the quantum channel by applying a programmatic rotation ($R_x$) of angle $\theta$ to Bob's logical qubit (qubit 2) prior to readout.\textsuperscript{1} The system executes a \texttt{teleport\_circuit} function, converts the full resulting state to a density matrix via \texttt{\_full\_state\_to\_density}, and performs a partial trace on qubits 0 and 1 to mathematically isolate Bob's resulting state.\textsuperscript{1} By converting this reduced density matrix back into a pure state, the software calculates the fidelity degradation compared to the original message state amplitudes ($\alpha, \beta$), computationally proving the ``no-cloning and monogamy'' principles of quantum mechanics without requiring access to a physical dilution refrigerator.\textsuperscript{1}

These highly abstract, NumPy-based \texttt{complex128} numerical arrays simulate perfect, idealized Bell states ($|\Phi^+\rangle = (|00\rangle + |11\rangle)/\sqrt{2}$).\textsuperscript{1} However, translating these idealized mathematical representations into physical hardware necessitates the exhaustive simulation of the cryogenic metamaterials that will physically host and route these states.

\section{Computational Modeling of Superconducting Magnetic Meta-Atoms}

The primary active physical components enabling dynamic microwave phase shifting and routing in the deep cryogenic regime are radio-frequency Superconducting Quantum Interference Device (rf-SQUID) arrays.\textsuperscript{1} Unlike a direct-current (DC) SQUID which contains two parallel Josephson junctions, an rf-SQUID consists of a single superconducting loop interrupted by exactly one highly nonlinear Josephson junction.\textsuperscript{1} When a dense, periodic array of these sub-wavelength devices is reactively coupled to a low-temperature superconductor (LTS) niobium microwave transmission line, the entire assembly functions as a flux-tunable effective medium.\textsuperscript{1} Simulating the emergent electromagnetic properties of this metamaterial requires multiscale computational approaches that bridge quantum Hamiltonian diagonalization and macroscopic finite-element electrodynamics.

\subsection{Simulating Quantum Energy Spectra with scqubits}

At the microscopic level, the behavior of an individual rf-SQUID meta-atom is governed fundamentally by quantum mechanics.\textsuperscript{4} In a traditional conductor, inductance is primarily a geometric property. However, in a superconductor, the inertia of the Cooper pairs carrying the supercurrent contributes an additional term known as kinetic inductance.\textsuperscript{1} The Josephson junction acts as a highly nonlinear inductor whose inductance value is strictly dependent on the gauge-invariant phase difference across the junction, which is in turn dictated by the external magnetic flux bias applied by the local Cryo-CMOS controller.\textsuperscript{1}

To extract the fundamental energy spectra, wavefunctions, and coherence times of these circuits, computational materials scientists heavily utilize \textit{scqubits}, an open-source Python package dedicated to the numerical analysis of superconducting circuits.\textsuperscript{6} The scqubits library allows computational physicists to programmatically define the Hamiltonian of the rf-SQUID and other common superconducting topologies (such as the transmon, fluxonium, $\cos(2\phi)$, and $0-\pi$ qubits).\textsuperscript{9}

The Hamiltonian for a generic flux-tunable superconducting loop can be mathematically modeled within the software as:
\begin{equation}
H = 4E_C n^2 - E_J \cos(\phi) + \frac{1}{2} E_L(\phi - \phi_{\text{ext}})^2
\end{equation}

In this computational model, $E_C$ represents the charging energy, $E_J$ is the Josephson energy, $E_L$ is the inductive energy of the loop, $\phi$ is the phase across the junction, $n$ is the conjugate charge operator, and $\phi_{\text{ext}}$ is the external magnetic flux bias.\textsuperscript{9} By truncating the theoretically infinite-dimensional Hilbert space to a computationally manageable localized basis, scqubits utilizes high-performance NumPy and SciPy backend eigensolvers---often accelerated via multi-core multiprocessing---to rapidly diagonalize the Hamiltonian.\textsuperscript{7}

This diagonalization provides the quantum energy level structure as a continuous function of the external flux bias. A critical computational insight derived from these programmatic parameter sweeps is the precise identification of the optimal flux bias points, commonly referred to as ``sweet spots''.\textsuperscript{12} At these coordinates, the software confirms that the non-linearity of the kinetic inductance is maximized while the first derivative of the transition frequency with respect to flux approaches zero, thereby minimizing the meta-atom's susceptibility to ambient $1/f$ flux noise.\textsuperscript{9} Furthermore, scqubits exposes a direct Application Programming Interface (API) to the Quantum Toolbox in Python (QuTiP). This integration enables the seamless simulation of composite Hilbert spaces, allowing researchers to computationally evaluate the time-domain evolution of the rf-SQUID meta-atom as it inductively couples to adjacent microwave cavity modes.\textsuperscript{7}

\subsection{Macroscopic Electrodynamics and 2D London Equation Solvers}

While scqubits effectively resolves the quantum dynamics and energy spectra of a single, isolated meta-atom, designing the Quantum ASIC requires simulating a massive 1,000-element metasurface array. Transitioning from a single component to a macroscopic array fundamentally alters the computational approach. Traditional 3D Finite Element Method (FEM) solvers often struggle with the extreme aspect ratios characteristic of thin-film superconductors.\textsuperscript{15} Attempting to mesh an ultra-thin superconducting layer within a large vacuum bounding box leads to combinatorial mesh explosions, consuming excessive RAM and resulting in intractable computation times.\textsuperscript{15}

To resolve this bottleneck, code-based material science leverages specialized 2D solvers such as \textit{SuperScreen}. SuperScreen is an open-source Python package, distributed under the MIT license, designed specifically to simulate the magnetic response of two-dimensional and quasi-2D superconductors.\textsuperscript{15} Rather than solving full 3D Maxwell equations, SuperScreen computationally solves the 2D London equation using an exceptionally efficient matrix inversion method.\textsuperscript{15}

The software algorithm computes the spatial distribution of Meissner screening currents and trapped magnetic flux given an applied, time-independent or quasi-DC magnetic field.\textsuperscript{15} The foundational physics equation solved by the software relates the 2D sheet current $K$ to the local magnetic vector potential $A$:
\begin{equation}
K(x,y) = -\frac{1}{\mu_0 \Lambda(x,y)} \left[ A(x,y) + \frac{\Phi_0}{2\pi} \nabla\theta(x,y) \right]
\end{equation}

In this computational framework, $\Lambda(x,y)$ represents the spatially inhomogeneous effective magnetic penetration depth, $\Phi_0$ is the magnetic flux quantum, $\mu_0$ is the vacuum permeability, and $\theta$ is the superconducting phase.\textsuperscript{15}

A profound advantage of SuperScreen is its programmatic ability to handle complex, irregular device geometries with spatially varying $\Lambda$ parameters.\textsuperscript{15} This is essential for modeling realistic rf-SQUID arrays where minor fabrication imperfections or variations in the Ginzburg-Landau critical temperature introduce localized disorder across the metasurface.\textsuperscript{5} By generating finite element data structures containing thousands of optimized vertices and triangles (e.g., a typical robust mesh utilizing 4,000 vertices and 8,000 triangles), SuperScreen rapidly calculates the self-inductance and mutual-inductance matrices of the entire 1,000-element array.\textsuperscript{15} These massive matrices provide the critical classical parameters required to compute the array's collective impedance and spatial phase delay characteristics, bridging the gap between quantum kinetics and macroscopic electrodynamics.\textsuperscript{15}

\subsection{High-Frequency Dynamics and Dissipative Breathers}

When the rf-SQUID array is actively driven by high-frequency microwave photons during a quantum routing operation, the weak magnetic dipole-dipole coupling between adjacent meta-atoms interacts with the intrinsic nonlinearity of the individual Josephson junctions.\textsuperscript{1} This complex interaction can lead to the spontaneous formation of dissipative discrete breathers---highly localized, spatially clustered states of electromagnetic excitation.\textsuperscript{1} When these breathers form, electromagnetic energy becomes violently trapped in a small subset of meta-atoms rather than propagating uniformly through the array, shattering the precise spatial phase gradients required to route quantum information across the bus.\textsuperscript{1}

Simulating these non-linear temporal dynamics and mitigating breather formation requires advanced simulation suites. Computational engineers utilize frequency-domain harmonic balance simulators, such as the Julia-based \textit{JosephsonCircuits.jl}, to calculate the non-linear small-signal response and scattering parameters of the highly coupled circuits via analytic Jacobians.\textsuperscript{21} For highly parallelized 3D evaluations, researchers deploy SQDMetal, a Python-based API that seamlessly integrates Qiskit Metal with the AWS Palace solver, Gmsh for meshing, and Paraview for visualization.\textsuperscript{23}

By extracting the Hamiltonian parameters via the Energy Participation Ratio (EPR) method and incorporating the kinetic inductance effects derived from scqubits, SQDMetal allows researchers to computationally identify the specific RF driving amplitudes and DC flux limits that induce breathers.\textsuperscript{5} Iterative simulations have proven that by computationally optimizing the inter-atomic coupling parameters and programming the Cryo-CMOS controllers to increase the amplitude of the applied RF flux, the system can force the array into a synchronized, coherent state, recovering the collective wavefront steering required by the Quantum ASIC.\textsuperscript{1}

\begin{table}[htbp]
\centering
\renewcommand{\arraystretch}{1.5}
\begin{tabularx}{\textwidth}{>{\raggedright\arraybackslash}p{3cm} >{\raggedright\arraybackslash}X >{\raggedright\arraybackslash}X >{\raggedright\arraybackslash}X}
\toprule
\textbf{Computational Tool} & \textbf{Primary Simulation Target} & \textbf{Underlying Mathematical Framework} & \textbf{Meta-Atom Application in Quantum ASIC} \\
\midrule
scqubits & Quantum energy spectra, coherence times & Hamiltonian diagonalization, Hilbert space truncation & Extracting $E_J$, $E_C$, and flux-dependent non-linear kinetic inductance.\textsuperscript{9} \\
SuperScreen & Meissner currents, self/mutual inductance matrices & 2D London equation, finite element matrix inversion & Mapping macroscopic magnetic screening and array characteristic impedance.\textsuperscript{15} \\
SQDMetal & RF S-parameters, device eigenmodes & Energy Participation Ratio (EPR), 3D FEM integration & Optimizing transmission line coupling and preventing localized dissipative breathers.\textsuperscript{23} \\
JosephsonCircuits.jl & Non-linear small-signal response & Harmonic balance method, Nodal analysis & Simulating parametric amplification and frequency conversion limits.\textsuperscript{21} \\
\bottomrule
\end{tabularx}
\caption{Simulation Stack for Superconducting Metamaterials}
\end{table}

\section{Phononic Metamaterials and Lithium Niobate Acoustic Dynamics}

Complementing the purely inductive response of superconducting rf-SQUIDs, the holographic quantum bus architecture also relies on advanced acoustic metamaterials to mediate microwave-to-phonon interactions.\textsuperscript{1} While microwave photons are excellent for fast communication across the bus, their long wavelengths make sub-millimeter spatial localization difficult. Phonons (quantized acoustic vibrations) operating at the same gigahertz frequencies possess wavelengths roughly five orders of magnitude smaller than their electromagnetic counterparts, allowing for incredibly dense, highly localized meta-atom arrays.\textsuperscript{1} Piezoelectric lithium niobate ($\text{LiNbO}_3$) Bulk Acoustic Wave (BAW) resonators have rapidly emerged as highly attractive solid-state meta-atoms for these applications due to their exceptionally strong electromechanical coupling coefficients ($k_t^2$).\textsuperscript{1}

\subsection{BeamProp and Acoustic Mode Propagation}

The physical integration of macroscopic lithium niobate BAW resonators into quantum circuits requires exhaustive material characterization.\textsuperscript{1} These resonators typically utilize a Z-cut crystalline orientation with a highly polished plano-convex geometry.\textsuperscript{1} The convex curvature acts as a critical acoustic lens, trapping the phonons in the center of the crystal and preventing them from radiating out into the physical supports, thereby drastically suppressing clamping losses.\textsuperscript{1}

Optimizing this geometry cannot rely on trial-and-error fabrication; it demands precise simulation of acoustic mode propagation. Computational physicists frequently adapt optical simulation libraries to study these mechanical modes, most notably leveraging tools related to the \textit{BeamProp} algorithm.\textsuperscript{26} BeamProp is a software package utilizing the Beam Propagation Method (BPM) based on robust implicit finite-difference schemes.\textsuperscript{26}

In the context of circuit Quantum Acoustodynamics (cQAD), the BPM algorithm represents the propagating acoustic wave as an angular spectrum of plane waves centered about the symmetry axis (the z-axis) of the High-overtone Bulk Acoustic Resonator (HBAR).\textsuperscript{26} By solving the paraxial wave equations, the Python-based scripts compute the spatial confinement of the two primary families of acoustic polarizations supported by the $\text{LiNbO}_3$ crystal symmetry: longitudinally polarized ``A-modes'' where bulk displacement occurs parallel to the Z-axis, and shear-wave polarized ``B-modes'' dominating the transverse X-Y plane.\textsuperscript{1} The computational output allows engineers to iteratively adjust the simulated radius of curvature to perfectly align the acoustic mode volume with the electric field maxima of the overlapping microwave split-post cavity, thereby maximizing the optomechanical coupling rate ($g_0$).\textsuperscript{1}

\subsection{The Landau-Rumer Attenuation Regime at Millikelvin Temperatures}

A defining parameter for acoustic metamaterial performance is the mechanical quality factor ($Q$). When operated in the high-frequency radio regime (1 to 100 MHz) and cooled to standard cryogenic temperatures (4 K), $\text{LiNbO}_3$ resonators exhibit extraordinary coherence.\textsuperscript{1} At 4 K, acoustic attenuation in the pristine crystal lattice is primarily governed by the Landau-Rumer microscopic model of phonon-phonon scattering.\textsuperscript{1} Because thermal phonon populations drop exponentially as the temperature decreases, the scattering cross-section plummets, theoretically leading to a massive spike in the mechanical $Q$ factor.\textsuperscript{1}

In computational acoustic simulations, the Landau-Rumer damping rate ($\gamma_{\text{LR}}$) can be approximated and coded into Python scripts as a strict temperature-dependent attenuation variable:
\begin{equation}
\gamma_{\text{LR}}^{-1} = \frac{\pi^5 \gamma^2 k_B^4 T^4}{15\rho v^5 h^3}
\end{equation}

In this formula, $\gamma$ represents the Gr\"uneisen parameter, $k_B$ is the Boltzmann constant, $T$ is the operational temperature, $\rho$ is the material density, $v$ is the phase velocity of sound in the material, and $h$ is Planck's constant.\textsuperscript{31} Advanced volumetric acoustic solvers, such as the Python-based \textit{OptimUS} library (which utilizes the boundary element method to solve the Helmholtz equation for harmonic wave propagation), can ingest these variables to predict precise wave decay and spatial transmission loss through the modeled piezoelectric layers over time.\textsuperscript{32} Based on these mathematical models, the $Q$ factor should continue to rise toward infinity as the temperature approaches absolute zero.

However, a profound material anomaly is observed when the simulation environment is computationally pushed down to 10 mK to match the operational environment of the superconducting qubits. Empirical data, replicated in advanced computational models, reveals that the quality factor violently degrades in the sub-Kelvin regime.\textsuperscript{1} The standard $T^4$ dependence of the Landau-Rumer equation catastrophically fails to predict this attenuation.\textsuperscript{1} The physical mechanism responsible for this anomalous millikelvin degradation is resonant absorption by microscopic Two-Level System (TLS) defects.\textsuperscript{1}

\subsection{The Sub-Kelvin Anomaly: Simulating Two-Level System (TLS) Defect Decoherence}

At 4 K, ambient thermal energy is high enough that microscopic structural impurities---such as dangling bonds, oxygen vacancies, and lattice dislocations within the $\text{LiNbO}_3$ crystal bulk or the amorphous aluminum oxide layers of the rf-SQUIDs---are constantly saturated and transitioning.\textsuperscript{1} This renders them largely transparent to the acoustic or microwave wave. However, as the temperature drops to 10 mK, these defects freeze into their absolute ground states, forming highly coherent Two-Level Systems (TLS).\textsuperscript{1}

When the Quantum ASIC meta-atom is excited, the propagating phonons or microwave photons are resonantly absorbed by these frozen defects. This interaction acts as a devastating source of $1/f$ dielectric noise, energy relaxation ($T_1$ decay), and phase decoherence ($T_2$ decay), severely limiting the circuit depth of any quantum algorithm.\textsuperscript{12} Understanding and mitigating this noise requires code-based material science to explicitly model the hardware as an open quantum system.

\subsection{Lindbladian Master Equations and QuTiP}

Because TLS defects represent a massive, complex environmental bath interacting with the primary quantum system, their dynamics cannot be modeled by simple unitary Schr\"odinger evolution.\textsuperscript{39} Instead, computational physicists model the hardware using the Lindblad master equation.\textsuperscript{39} The time evolution of the system's density matrix $\rho(t)$ is computationally tracked as:
\begin{equation}
\dot{\rho} = -\frac{i}{\hbar}[H_{\text{sys}} + H_{\text{bath}} + H_{\text{int}}, \rho] + \sum_k \left( L_k \rho L_k^\dagger - \frac{1}{2}\{L_k^\dagger L_k, \rho\} \right)
\end{equation}

In this programmatic model, $H_{\text{sys}}$ is the qubit or meta-atom Hamiltonian, $H_{\text{bath}}$ represents the ensemble of TLS defects, $H_{\text{int}}$ represents the dipole-coupling between the system and the defects, and $L_k$ are the collapse (jump) operators describing irreversible energy loss and dephasing to the macroscopic environment.\textsuperscript{39}

Executing a full Lindbladian simulation for a physical material containing millions of distributed TLS defects is computationally intractable on classical hardware.\textsuperscript{39} Therefore, advanced Python simulation frameworks invoke strict statistical truncation based on the phenomenological standard tunneling model.\textsuperscript{38} The simulation code mathematically generates a spatial distribution of, for example, one million TLSs across the surface of the simulated qubit.\textsuperscript{39} The algorithm calculates the localized electric field distributions (using FEM solvers) and identifies the specific subset of TLSs that dominate the interaction---typically the 200 most strongly coupled defects located near the Josephson junctions or within areas of maximum acoustic strain.\textsuperscript{39} The continuous Hilbert space is truncated, and the Lindbladian is solved exclusively for this strongly coupled sub-ensemble.\textsuperscript{39}

Open-source libraries like the Quantum Toolbox in Python (QuTiP) are highly optimized for this precise task.\textsuperscript{9} QuTiP provides built-in \texttt{DecoherenceNoise} classes and predefined collapse operators to numerically integrate the master equation over time, allowing researchers to extract the bi-exponential decay curves that perfectly match the $T_1$ and $T_2$ relaxation times observed in empirical 10 mK dilution refrigerator experiments.\textsuperscript{14}

\textbf{Computational Insights and TLS Mitigation:} A profound higher-order insight derived directly from these simulations is that strongly coupled TLSs located far from the Josephson junction can still significantly degrade coherence if they happen to reside in areas of highly concentrated electric fields.\textsuperscript{39} The computational model proves that simply moving defects away from the junction is insufficient; the global field geometry must be addressed.\textsuperscript{39}

Consequently, code-based mitigation strategies have emerged. By utilizing \textit{QwaveMPS} (an open-source Python library for simulating one-dimensional quantum many-body waveguide systems using Matrix Product States), researchers can efficiently simulate the decay of TLSs in non-Markovian waveguide QED regimes.\textsuperscript{43} These Matrix Product State simulations have validated that engineering a phononic bandgap in the underlying silicon-on-insulator substrate can completely suppress all microwave-frequency phonons around the operating frequency of the transmon or meta-atom.\textsuperscript{44} By actively tuning the TLS frequencies out of the spectral vicinity of the qubits via applied local DC electric fields (simulated in tandem with the open-system dynamics), the system effectively forces the parasitic defects into an acoustic vacuum.\textsuperscript{44} This computational insight, when translated to physical hardware, has resulted in profound performance enhancements, elevating relaxation times by two orders of magnitude and suppressing TLS-induced performance outliers by a factor of four.\textsuperscript{44}

\section{Algorithmic Topology Optimization: QAOA and QUBO Implementations}

With the physical dynamics of the cryogenic meta-atoms and their parasitic defects accurately modeled and optimized, the next massive computational challenge is the high-level control of the Quantum ASIC.\textsuperscript{1} The fundamental purpose of the architecture is to establish a software-defined quantum bus to solve the physical wiring crisis.\textsuperscript{1} However, achieving dynamic, all-to-all connectivity without permanent physical traces requires a highly sophisticated algorithmic framework capable of mapping logical quantum circuits onto the physical coordinates of the metasurface.\textsuperscript{1}

Because physical qubits on a solid-state chip can only interact with their immediate nearest neighbors, bridging non-adjacent logical qubits requires executing transient SWAP gates or dynamically programming the holographic metasurface to create a transient electromagnetic channel.\textsuperscript{1} Determining the absolute most efficient spatial assignment of logical roles to physical qubits to minimize the total gate depth and SWAP count is a computational problem that scales exponentially, rapidly becoming NP-hard.\textsuperscript{48}

\subsection{The QUBO Formulation and the QAOA Solver}

To solve this, the quantum routing challenge is mathematically formulated as a Quadratic Unconstrained Binary Optimization (QUBO) problem.\textsuperscript{1} In a QUBO model, binary decision variables (representing specific qubit assignments or edge routings on the metasurface) are mapped to a rigorous cost function.\textsuperscript{51} This cost function is mathematically designed to heavily penalize inefficient wiring distances, excessive SWAP gate insertion, and deep circuit depths that would otherwise allow the quantum state to succumb to TLS decoherence before computation is complete.\textsuperscript{1}

The optimization of this QUBO is executed using the Quantum Approximate Optimization Algorithm (QAOA), a hybrid quantum-classical variational algorithm highly suited for near-term noisy superconducting hardware.\textsuperscript{1} In this variational execution, the QAOA algorithm navigates the massive $2^N$-dimensional Hilbert space of all possible routing configurations across the metasurface.\textsuperscript{1}

The execution begins by initializing the quantum hardware in the maximally superimposed ground state of a Mixing Hamiltonian ($H_M$), typically defined as a uniform transverse-field Ising model composed entirely of Pauli-X operators:
\begin{equation}
H_M = \sum_{i=1}^N \sigma_i^x
\end{equation}

Physically, this initialization implies that every possible routing configuration on the metasurface exists simultaneously in quantum superposition.\textsuperscript{1} The algorithm then undergoes unitary evolution, continuously alternating between the Mixing Hamiltonian (which allows the system to quantum-tunnel out of sub-optimal local minima energy traps) and a Problem Hamiltonian ($H_P$, which encodes the QUBO cost function into a quantum-computable format via commuting Pauli-Z operators).\textsuperscript{1}

\subsection{Qiskit and OpenQAOA Code Implementations}

The practical software implementation of this routing protocol relies heavily on IBM's Qiskit SDK and the \textit{OpenQAOA} Python library.\textsuperscript{47} In software, the initial logical quantum circuit is first modeled as a classical graph using network libraries (e.g., \texttt{rx.PyGraph}), where nodes represent the logical qubits and edges represent the required entanglement interactions dictated by the quantum algorithm.\textsuperscript{52}

Using the \texttt{qiskit\_optimization} module, this graph is ingested and converted into a \texttt{QuadraticProgram} via \texttt{problem.binary\_var()} assignments.\textsuperscript{52} The software then automatically translates this classical optimization problem into an Ising Hamiltonian using the built-in \texttt{to\_ising()} method, generating the sparse Pauli operations that the quantum hardware can natively process.\textsuperscript{51}

Advanced transpilation passes and custom Python routing functions (e.g., \texttt{my\_custom\_routing\_function}) are coded to define \texttt{swap\_mask} logic.\textsuperscript{47} A critical computational optimization realized in recent Qiskit deployments is the mathematical proof that multi-qubit CPHASE gates in QAOA circuits are perfectly commutative.\textsuperscript{48} By exploiting this commutativity, integer programming algorithms within the transpiler can aggressively reorder the gate execution sequence, enabling massive parallel execution and significantly reducing the total required SWAP gate count.\textsuperscript{48}

Through an iterative quantum-classical loop managed by Qiskit primitives (like the Sampler), a classical optimizer (such as COBYLA) adjusts the variational evolution angles ($\beta, \gamma$) to maximize constructive quantum interference around the absolute most efficient routing layout.\textsuperscript{1} Empirical data from IBM cloud executions (specifically benchmarking on the 127-qubit \texttt{ibm\_fez} hardware) demonstrate that this hybrid algorithm successfully collapses the wavefunction to an optimal logical-to-physical mapping vector with exceptional objective penalty scores, routinely achieving values as low as 4.0.\textsuperscript{1}

\begin{table}[htbp]
\centering
\renewcommand{\arraystretch}{1.5}
\begin{tabularx}{\textwidth}{>{\raggedright\arraybackslash}p{3cm} >{\raggedright\arraybackslash}X >{\raggedright\arraybackslash}X >{\raggedright\arraybackslash}X}
\toprule
\textbf{Algorithmic Component} & \textbf{Python Software Implementation} & \textbf{Mathematical Representation} & \textbf{Function in Quantum ASIC Architecture} \\
\midrule
Problem Formulation & Qiskit QuadraticProgram & QUBO Matrix Formulation & Mathematically penalizes inefficient layout topology and long trace distances.\textsuperscript{53} \\
Problem Hamiltonian & Qiskit \texttt{to\_ising()} method & Pauli-Z operators & Encodes the classical cost function into a quantum-computable energy landscape.\textsuperscript{51} \\
Mixing Hamiltonian & Qiskit Circuit Ansatz & Pauli-X operators ($H_M = \sum \sigma^x$) & Drives global state transitions to avoid trapping the optimizer in local minima.\textsuperscript{1} \\
Optimization Loop & OpenQAOA / COBYLA optimizer & Unitary Evolution ($\beta, \gamma$) & Iteratively collapses the wavefunction to reveal the optimal spatial routing vector.\textsuperscript{47} \\
\bottomrule
\end{tabularx}
\caption{QAOA Implementation Stack for Topology Routing}
\end{table}

\section{Deep Neural Network Surrogate Modeling for Real-Time Holography}

Once the QAOA solver successfully extracts the optimal low-dimensional routing topology vector (e.g., coordinates explicitly mapping Qubit 0 to Qubit 2), a massive computational hurdle remains.\textsuperscript{1} This low-dimensional data must be instantly translated into high-dimensional, continuous electromagnetic phase commands for the 1,000 individual meta-atoms comprising the holographic metasurface.\textsuperscript{1}

This translation poses the most severe computational bottleneck in the entire Quantum ASIC architecture. Traditional computational electromagnetics relies on solving the full-wave Maxwell equations using techniques such as the Finite-Difference Time-Domain (FDTD) method, Finite Element Method (FEM), or rigorous coupled-wave analysis.\textsuperscript{1} Calculating the complex scattering matrices and near-field mutual coupling for an aperiodic array of 1,000 mutually interacting rf-SQUIDs or BAW resonators takes hours or days on classical supercomputers.\textsuperscript{1} This extreme latency is fundamentally incompatible with the microsecond-scale coherence times ($T_2$) of superconducting qubits; the fragile quantum state would completely decohere long before the classical FDTD solver could compute the required routing phase profile.\textsuperscript{1}

\subsection{PyTorch Architecture and Forward/Inverse Design}

To achieve real-time dynamic holography and bypass the FDTD bottleneck entirely, code-based material science leverages Deep Neural Network (DNN) surrogate modeling.\textsuperscript{1} Implemented using modern deep learning frameworks such as PyTorch (e.g., version 1.1.0 or 1.8.0) or TensorFlow, these neural networks act as highly efficient mathematical surrogates that instantly infer the complex, cross-coupling electromagnetic interactions between meta-atoms without solving Maxwell's equations.\textsuperscript{1}

The implementation workflow involves two distinct phases: forward modeling and inverse design.\textsuperscript{56}

\textbf{Forward Prediction Network (FPN):} Initially, a massive dataset is generated using standard FDTD/FEM solvers (such as COMSOL or SQDMetal), correlating specific meta-atom geometric configurations and flux bias states to their resulting S-parameters (transmission amplitude, phase shift, and near-field radiation patterns).\textsuperscript{60} Extensive numerical experiments reveal that a relatively sparse dataset comprising approximately 4,000 samples is adequate to establish a robust DNN model capable of greater than 90\% accuracy.\textsuperscript{62} A Convolutional Neural Network (CNN) is trained on this dataset using an optimizer like Adam. The loss function during training (typically Mean Squared Error, MSE) rigorously minimizes the mathematical difference between the FEM-simulated S-parameters and the network's internal predictions.\textsuperscript{3} Once trained, the FPN can predict the electromagnetic response of an arbitrary metasurface configuration in milliseconds.\textsuperscript{59}

\textbf{Inverse Design Network:} Once the forward model is validated, the architecture is flipped. The optimal routing vector from the QAOA algorithm---represented as a normalized topology feature vector $\mathbf{T}_{\text{target}} \in \mathbb{R}^d$, where $d=10$ features representing spatial coordinates and required coupling strengths---is ingested into the input layer of a tandem network (often integrating Multi-Layer Perceptrons or Variational Autoencoders).\textsuperscript{1} The hidden layers ($L_1, L_2, \dots, L_n$) utilize non-linear activation functions like Rectified Linear Units (ReLU, where $\sigma(x) = \max(0, x)$) to construct a highly complex latent representation of the desired scattering matrix:
\begin{equation}
\mathbf{H}_l = \sigma(\mathbf{W}_l \mathbf{H}_{l-1} + \mathbf{b}_l)
\end{equation}
where $\mathbf{W}_l$ represents the weight matrix and $\mathbf{b}_l$ the bias vector for layer $l$.\textsuperscript{1}

\textbf{Phase Synthesis and Thermodynamic Validation:} The final output layer of the PyTorch inverse model utilizes a Sigmoid activation function to map the latent representation to a continuous tensor representing the physical phase shifts required for all $M$ meta-atoms simultaneously:
\begin{equation}
\mathbf{\Phi}_{\text{pred}} = \text{Sigmoid}(\mathbf{W}_{\text{out}} \mathbf{H}_n + \mathbf{b}_{\text{out}}) \times 2\pi
\end{equation}
This equation instantly outputs an exact phase array $\mathbf{\Phi}_{\text{pred}} \in [0, 2\pi]^M$ (e.g., $M = 1000$ sub-wavelength elements).\textsuperscript{1}

A profound third-order insight emerged during the empirical benchmarking of this generative PyTorch model. When tasked with synthesizing the 1,000-element phase profile required to route the microwave photons, the DNN did not output a broad, uniform distribution of phases spanning the full $0$ to $2\pi$ dynamic range. Instead, the network generated phase shifts that were strictly concentrated in a highly narrow band between 3.033 and 3.284 radians, with a precise mean phase of 3.142 radians ($\approx \pi$).\textsuperscript{1}

This computational output represents a thermodynamic triumph. A uniform $\pi$ phase shift applied across an entire surface effectively acts as a perfect flat inversion mirror. By computationally applying only subtle, sub-radian perturbations (e.g., 3.03 or 3.28 rad) around this $\pi$-radian baseline, the metasurface leverages slight geometric phase deviations to create the precise constructive and destructive interference patterns required to focus the beam on the target qubit.\textsuperscript{1}

Through unsupervised learning, the deep neural network effectively deduced the thermodynamic limits of the hardware. To physically actuate these phase shifts in the dilution refrigerator, the heterogeneously integrated 28nm FD-SOI Cryo-CMOS controller utilizes Charge-Lock Fast-Gate (CLFG) cells that toggle trapped charge on floating capacitors.\textsuperscript{1} Because the DNN restricts the required phase tuning to such a narrow margin, the Cryo-CMOS cells are not required to inject massive amounts of voltage to drive broad phase sweeps.\textsuperscript{1} This computational constraint explicitly preserves the ultra-low 18 nanowatt (nW) active dissipation limit per cell, ensuring the total 1,000-channel controller dissipates merely 18 microwatts ($\mu$W)---well below the failure threshold of the 10 mK dilution refrigerator.\textsuperscript{1} Without the DNN's mathematically constrained phase generation, the Quantum ASIC would instantly crash the cryostat via thermal occlusion.

\section{Macroscopic Extensions: Code-Driven Quantum SATCOM and Quantum Radar}

The exact same computational physics frameworks---the rf-SQUID Hamiltonian solvers, the 2D London equation evaluations, the Lindbladian open-system TLS models, and the PyTorch inverse design phase synthesizers---scale directly from the internal sub-millimeter quantum bus to massive, macroscopic phased-array apertures used for environmental sensing and global Quantum Satellite Communications (SATCOM).\textsuperscript{1}

Global optical Quantum Key Distribution (QKD) currently suffers critically from Mie scattering, where near-infrared photons are deflected out of the communication channel by atmospheric aerosols of comparable physical size (e.g., clouds and fog).\textsuperscript{1} Code-based simulations of microwave propagation shift the communication protocol to the Rayleigh Scattering Approximation ($\gamma_c(f,T) = K_l(f,T)M$).\textsuperscript{1} Under Rayleigh dynamics, centimeter-wavelength microwaves are vastly larger than atmospheric aerosols, allowing the photons to effectively diffract around suspended water particles, yielding a fundamentally weather-resilient channel.\textsuperscript{1}

However, ambient thermal noise at microwave frequencies in a 300 K Earth environment generates a massive population of thermal microwave photons that instantly drown out fragile continuous-variable (CV) quantum states.\textsuperscript{1} To compute a solution, researchers integrate thermodynamic models with the metasurface's routing algorithms to execute ``radiative cooling'' protocols.\textsuperscript{1} By highly directing the active metasurface aperture and computationally overcoupling the microwave communication channel to a 10 mK cold load (or the 2.7 K vacuum of space), the system actively drains thermal noise out of the transmission line.\textsuperscript{1} Simulation code confirms this protocol successfully manipulates Bose-Einstein statistics to drive the thermal photon occupation number down to an exceptional 0.06, permitting unconditionally secure CV state teleportation with state transfer fidelities of 58.5\% even when the metasurface is operating in 4 K ambient conditions.\textsuperscript{1}

Furthermore, extending these numerical models to incorporate two-mode squeezed vacuum (TMSV) states enables the deployment of Quantum Illumination, commonly referred to as Quantum Radar.\textsuperscript{1} The code simulates the generation of highly correlated signal and idler photons. The signal photon is transmitted into the noisy environment to scan for targets, while the idler is retained locally in a cryogenic memory.\textsuperscript{1} Although environmental decoherence destroys the pure quantum entanglement during transit, mathematical simulations of non-classical correlation metrics---specifically quantum discord---prove that the joint measurement of the scattered return signal against the locally retained idler effectively doubles the quantum Chernoff error exponent compared to the best possible classical heterodyne receiver.\textsuperscript{1} This algorithmic advantage elevates detection precision to the theoretical Heisenberg limit ($1/M$), empowering the microwave metasurface aperture to reject urban multipath clutter that is up to 20 dB stronger than true target echoes, achieving a 1.5x target range extension in dense fog simulations.\textsuperscript{1}

\section{Hardware-Software Co-Design Ecosystems: Identifying Potential Partners}

The successful execution of this highly integrated, code-based material science paradigm requires an intense geographic and intellectual concentration. The translation of PyTorch surrogate models, QAOA topologies, and scqubits Hamiltonians into physical, space-qualified 18-nanowatt controllers and lithium niobate BAWs cannot occur in isolated academic silos. Regions with dense clusters of computational chemistry, advanced nanofabrication, and algorithmic design serve as critical accelerators for the deployment of these technologies.

The academic and industrial ecosystem in San Diego, California, serves as an excellent model for this vital convergence. At the computational tier, leveraging facilities analogous to the San Diego Supercomputer Center (SDSC) at the University of California, San Diego (UCSD) would allow for the development of advanced GPU-accelerated molecular simulation software capable of performing the high-throughput Quantum Mechanics/Molecular Mechanics (QM/MM) calculations required to analyze dielectric impurities at the atomic scale. These frameworks could feed directly into our Lindbladian TLS models used to mitigate decoherence across the quantum bus.

Simultaneously, engaging with advanced characterization labs, such as UCSD's Allen Lab, would provide the essential empirical validation loop required to verify Python simulations. By employing low-temperature Scanning Microwave Microscopy (SMM) operating at 50 mK, these types of researchers directly visualize the dissipation, screening responses, and edge magnetoplasmons of topological states and metasurfaces under high-frequency electromagnetic fields. This physical near-field data could be ingested back into our PyTorch CNN forward prediction networks, continuously refining the accuracy of the S-parameter loss functions.

At the commercial manufacturing tier, partnering with specialized startups will be necessary to bridge the immense gap between simulated code and deployable hardware. Companies similar to Monarch Quantum, which focus aggressively on the robotic assembly and AI-driven machine-learning alignment of integrated photonics, represent the ideal profile for such partnerships. Supported by massive federal mandates---such as providing systems for NASA missions---engaging with such an ecosystem ensures that the theoretical phase arrays synthesized by our DNN algorithms can be physically instantiated into functional neutral-atom quantum sensors and SATCOM hardware capable of surviving severe launch loads and orbital vibration environments.

\section*{Cost-Effective Computational Infrastructure Alternatives}

While the computational requirements for simulating cryogenic quantum metamaterials---including 3D FEM solvers, QAOA optimization loops, and Lindbladian master equations---are substantial, several cost-effective alternatives can be utilized to run these workloads efficiently:

\begin{itemize}
    \item \textbf{Specialized GPU Cloud Providers:} Training the PyTorch Deep Neural Network surrogate models requires significant GPU acceleration. Instead of relying on traditional major cloud providers (AWS, GCP, Azure), utilizing specialized GPU clouds such as Lambda Labs, RunPod, or CoreWeave can provide access to high-tier compute (e.g., NVIDIA A100s or H100s) at a fraction of the hourly cost.
    \item \textbf{Spot and Preemptible Instances:} For highly parallelizable and fault-tolerant workloads---such as \textit{scqubits} parameter sweeps, QUBO objective evaluations, and QuTiP Monte Carlo trajectory simulations---leveraging Spot Instances (AWS) or Preemptible VMs (Google Cloud) offers massive computational power at up to 90\% discounts compared to on-demand pricing.
    \item \textbf{Academic Supercomputing Networks:} For foundational research, partnering with academic institutions can grant access to federated computing grids. The ACCESS (Advanced Cyberinfrastructure Coordination Ecosystem: Services \& Support) program or the Open Science Grid (OSG) provide virtually free, high-throughput computing allocations for researchers in the United States.
    \item \textbf{Cloud-Native Quantum Simulators:} Before deploying algorithms on expensive physical Quantum Processing Units (QPUs), extensive QAOA routing and state vector simulations can be run locally or on free-tier cloud simulators (such as IBM Quantum's \textit{qiskit-aer}) to validate the topology optimization algorithms at no cost.
\end{itemize}

\section{Strategic Conclusions}

The scalability crisis threatening the future of fault-tolerant quantum computing and secure global communication networks is fundamentally a materials and thermodynamics problem. Static coaxial architectures and optical atmospheric limitations cannot be overcome by traditional iterative engineering; the strict micro-watt cooling capacities of the 10 mK cryogenic environment absolutely prohibit it. The paradigm shift toward a holographic Quantum ASIC and macroscopic quantum SATCOM apertures provides the necessary physical solution, but it is the underlying code-based material science that provides the capability to design it.

Through the rigorous integration of open-source Python libraries, computational physicists can unravel the deep quantum dynamics of rf-SQUID meta-atoms (\textit{scqubits}), execute massive matrix inversions to map 2D Meissner screening (\textit{SuperScreen}), and precisely model the deeply non-linear acoustic attenuation of lithium niobate BAWs outside of the traditional Landau-Rumer regime (\textit{BeamProp, COMSOL}). By combining these solid-state simulations with open-system Lindbladian dynamics (\textit{QuTiP}), the devastating impacts of Two-Level System defects can be computationally isolated and effectively mitigated via active phononic bandgap tuning.

The overarching control of these advanced materials is entirely algorithmic. Formulating the ``wiring crisis'' as a QUBO problem allows QAOA solvers on cloud-based quantum hardware (\textit{Qiskit, OpenQAOA}) to instantly deduce optimal topological mappings. By feeding this routing data into highly optimized PyTorch Multi-Layer Perceptrons, the architecture completely bypasses hours of classical Maxwell equation solving. The neural networks instantly infer the sub-wavelength cross-coupling of the metasurface, generating precise, continuous phase shifts in milliseconds.

Crucially, by computationally constraining these phase shifts around a $\pi$-radian baseline, the deep learning model inadvertently enforces the strict thermodynamic limits of the 28nm FD-SOI Cryo-CMOS controllers, preserving the ultra-low 18 nanowatt thermal budget required for operation. Supported by vertically integrated research hubs spanning from GPU-accelerated molecular dynamics to space-qualified hybrid packaging, this comprehensive software-to-hardware pipeline successfully establishes the foundational engineering framework for realizing a distributed, highly scalable, and unconditionally weather-resilient quantum internet.

\vspace{0.5em}
\noindent\textbf{Related reports.} The Markdown whitepaper \textit{Holographic Metasurfaces as a Scalable Control Layer for Solid-State Quantum Entanglement and Secure SATCOM} and the long-form LaTeX report \textit{Cryogenic Metamaterial Architectures for Solid-State Quantum Routing and SATCOM} are companion documents in the QASIC Engineering-as-Code repository; see the repository README for links.

\section*{Works Cited}
\begin{enumerate}
    \item Cryogenic Metamaterial Architectures Quantum SATCOM
    \item Deep neural networks for the evaluation and design of photonic devices - UMBC (Nat. Reviews 2020)
    \item Deep inverse photonic design - OSTI
    \item Superconducting Quantum Interference Device (SQUID) - UBC Physics \& Astronomy
    \item Coherent oscillations of driven rf SQUID metamaterials - Institute of Radio-engineering and Electronics of RAS
    \item Superconducting Qubits in Python - scQubits documentation
    \item Scqubits: a Python package for superconducting qubits - ResearchGate
    \item scqubits: a Python package for superconducting qubits - arXiv
    \item Scqubits: a Python package for superconducting qubits - Quantum Journal
    \item A quantum engineer's guide to superconducting qubits - Applied Physics Reviews
    \item Studying phonon coherence with a quantum sensor - PMC - NIH
    \item 1/f Noise from Condensed Matter Physics to Quantum Technologies - Agenda INFN
    \item qutip-qip/doc/source/qip-processor.rst - GitHub
    \item SuperScreen: An open-source package for simulating the magnetic response of two-dimensional superconducting devices - OSTI.GOV
    \item Pyton TDGL - Scribd
    \item pyTDGL: Time-dependent Ginzburg-Landau in Python - SSRN
    \item Magnetic-field and current-density distributions in thin-film superconducting rings and disks
    \item Wide Band Shear Bulk Acoustic Resonators Using Crystalline Y-cut Lithium Niobate thin films with spurious suppression
    \item rf SQUID Metamaterials: A Rich Nonlinear Setting for Applications - Superconductivity News Forum
    \item JosephsonCircuits.jl: Frequency domain, multi-tone harmonic balance - GitHub
    \item QC-Hardware-How-To - GitHub
    \item Open-Source Highly Parallel Electromagnetic Simulations for Superconducting Circuits
    \item SQDMetal: Tools to aid in simulating and fabricating superconducting quantum devices - GitHub
    \item Development of Thin Film Lithium Niobate Shear Bulk Acoustic Wave Resonators
    \item Characterization of Lithium Niobate for High-Overtone Bulk Acoustic Wave Devices - ETH Z\"urich
    \item BeamPROP - micro.fel.cvut.cz
    \item Quartz phononic crystal resonators for hybrid acoustic quantum memories
    \item Low Temperature Properties of Low-Loss Macroscopic Lithium Niobate Bulk Acoustic Wave Resonators - arXiv
    \item Loss mechanisms in high-coherence multimode mechanical resonators coupled to superconducting circuits - arXiv
    \item UKAN+ Webinar: OptimUS: an open-source Python library for 3D acoustic wave propagation
    \item A mechanical quantum memory for microwave photons - arXiv
    \item Evidence of Memory Effects in the Dynamics of Two-Level System Defect Ensembles Using Broadband, Cryogenic Transient Dielectric Spectroscopy - arXiv
    \item Locating Two-Level Systems in a Superconducting Xmon Qubit - MDPI
    \item In situ scanning gate imaging of individual quantum two-level system defects in live superconducting circuits - PMC
    \item Simulating noise on a quantum processor: interactions between a qubit and resonant two-level system bath - OSTI
    \item Pulse-level noisy quantum circuits with QuTiP - Chalmers
    \item Mapping Two-Level-Systems On Transmon Qubits Reveals Surface Defects Primarily Reside On Josephson Junction Leads - Quantum Zeitgeist
    \item QwaveMPS: An efficient open-source Python package for simulating non-Markovian waveguide-QED using matrix product states - arXiv
    \item Phonon engineering of atomic-scale defects in superconducting quantum circuits - PMC
    \item Phonon-protected superconducting qubits - EECS Berkeley
    \item Scalable and Site-Specific Frequency Tuning of Two-Level System Defects in Superconducting Qubit Arrays
    \item Qubit Routing for QAOA - OpenQAOA
    \item Portfolio Optimization using QUBO and Qiskit
    \item Implementing the Quantum Approximate Optimization Algorithms for QUBO problems Across Quantum Hardware Platforms - arXiv
    \item Grover Optimizer - Qiskit Optimization
    \item Quantum-Assisted Vehicle Routing: Realizing QAOA-based Approach on Gate-Based Quantum Computer - arXiv
    \item Quantum approximate optimization algorithm (QAOA) - IBM
    \item Mapping a QUBO onto an observable for use with Qiskit QAOA
    \item Metanet - Stanford University
    \item Exploring AI in metasurface structures with forward and inverse design - PMC
    \item Inverse design of polarization conversion metasurfaces by deep neural networks
    \item Deep Learning-Assisted Design for High-Q-Value Dielectric Metasurface Structures - MDPI / PMC
    \item Deep Learning Enables Multifunctional Metasurfaces Design with Mutual Coupling Estimation - IEEE Xplore
    \item Inverse Design of Metasurfaces Using Machine Learning Techniques
    \item Prediction Enhancement of Metasurface Absorber Design Using Adaptive Cascaded Deep Learning (ACDL) Model - MDPI
    \item Inverse-design-of-metasurfaces - GitHub
    \item Quantum Foundry - SCIL
    \item A Quantum Leap for Molecular Simulations on GPUs - UC San Diego
    \item Research - Allen Lab
    \item Integrated Photonics Company Launches in San Diego
    \item Monarch Quantum Selected to Deliver Integrated Photonics Quantum Light Engines for NASA's First Planned Space-Based Quantum Gravity Gradiometer
\end{enumerate}

\end{document}
